\section*{Sets and Indices}

\begin{itemize}
    \item $W$: set of workshops, indexed by $w$. In this instance,
    \[
    W=\{A,B,C,D\}.
    \]
    \item $C$: set of components, indexed by $c$. In this instance,
    \[
    C=\{1,2,3\}.
    \]
\end{itemize}

\section*{Parameters}

\begin{itemize}
    \item $H_w$: available production capacity of workshop $w$ in hours (code: \texttt{capacities[w]}).
    \[
    H_A=100,\quad H_B=150,\quad H_C=80,\quad H_D=200.
    \]
    \item $r_{wc}$: production rate of component $c$ at workshop $w$ in units per hour (code: \texttt{rates[(w, c)]}).
    
    For this instance,
    \[
    \begin{array}{c|ccc}
    w & r_{w1} & r_{w2} & r_{w3} \\
    \hline
    A & 10 & 15 & 5 \\
    B & 15 & 10 & 5 \\
    C & 20 & 5 & 10 \\
    D & 10 & 15 & 20
    \end{array}
    \]
\end{itemize}

\section*{Decision Variables}

\begin{itemize}
    \item $x_{wc} \ge 0$: hours allocated by workshop $w$ to production of component $c$ (code: \texttt{x[(w, c)]}).
    \item $z \ge 0$: number of completed products (code: \texttt{z}).
\end{itemize}

\section*{Objective}

Maximize the number of completed products:
\[
\max \; z.
\]

\section*{Constraints}

\begin{align}
\sum_{c \in C} x_{wc} &\le H_w
&& \forall w \in W
\label{cap}\\
z &\le \sum_{w \in W} r_{wc}\, x_{wc}
&& \forall c \in C
\label{mincomp}\\
x_{wc} &\ge 0
&& \forall w \in W,\; c \in C
\label{nonnegx}\\
z &\ge 0.
\label{nonnegz}
\end{align}

Constraint \eqref{cap} enforces each workshop's available hours.  
Constraint \eqref{mincomp} models completed products as the minimum of total units produced across the three components, by requiring $z$ to be no larger than the total production of each component.

\section*{Notes on Implementation}

The implemented model in \texttt{current\_heuristic.py} is a linear program solved with PuLP using the Gurobi command-line solver interface (\texttt{GUROBI\_CMD}). The code does \emph{not} impose integrality on either the hour-allocation variables $x_{wc}$ or the completed-product variable $z$; all are continuous and nonnegative.

The model does not include separate component-balance equations beyond the inequalities $z \le \sum_{w \in W} r_{wc}x_{wc}$. Thus, exactly as implemented, the objective maximizes a continuous variable $z$ representing the minimum available component quantity across components.

For this specific instance, the LP is therefore:
\[
\max z
\]
subject to
\begin{align*}
x_{A1}+x_{A2}+x_{A3} &\le 100,\\
x_{B1}+x_{B2}+x_{B3} &\le 150,\\
x_{C1}+x_{C2}+x_{C3} &\le 80,\\
x_{D1}+x_{D2}+x_{D3} &\le 200,\\
z &\le 10x_{A1}+15x_{B1}+20x_{C1}+10x_{D1},\\
z &\le 15x_{A2}+10x_{B2}+5x_{C2}+15x_{D2},\\
z &\le 5x_{A3}+5x_{B3}+10x_{C3}+20x_{D3},\\
x_{wc} &\ge 0 \qquad \forall w\in W,\; c\in C,\\
z &\ge 0.
\end{align*}
